%!TEX program = lualatex
\documentclass[ngerman,headheight=70pt]{scrartcl}
\usepackage[ngerman]{babel}
\usepackage{blindtext}
\usepackage[utf8]{luainputenc}
\usepackage{scrlayer-scrpage}
\usepackage[T1]{fontenc}
\usepackage{fontspec}
\usepackage{multicol}
\usepackage{enumitem}
\usepackage[german=quotes]{csquotes}
\usepackage[bottom=2cm,footskip=8mm]{geometry}
\setmainfont{Source Sans Pro}
\MakeOuterQuote{"}

\setlength{\parindent}{0pt}
\setlength{\parskip}{6pt}

\pagestyle{plain.scrheadings}
\thispagestyle{scrheadings}
\lohead{\vspace*{-1.5cm}\includegraphics[scale=0.4]{../images/stupa_logo_rot.pdf}}
\rohead{\\Präsidium des\\Studierendenparlaments}

\lofoot[Universität Hamburg · Präsidium des Studierendenparlaments\\\hrule Von-Melle-Park 5 · D-20146 Hamburg]{Universität Hamburg · Präsidium des Studierendenparlaments\\\hrule Von-Melle-Park 5 · D-20146 Hamburg}
\cofoot[]{}
\cohead[\thepage]{}

\renewcommand{\thesubsection}{TOP \arabic{subsection}}
\renewcommand{\thesubsubsection}{\arabic{subsubsection}.}
\setcounter{subsection}{-1}

\begin{document}
    UHH · StuPa-Präsidium · Von-Melle-Park 5 · D-20146 Hamburg

    %%% PAGE ONE (EDIT LATER) %%%%
    % Hier den Titel des Protokolls eintragen
    \section*{Protokoll der 3. Sitzung des Studierendenparlaments vom 26. Mai 2016}

    % Protokollanten eintragen
    \textbf{Protokoll: Jim Martens}\\
    \textbf{Ort: S08 in der ehemaligen HWP}\\
    \textbf{Beginn: 18.39 Uhr}\\
    \textbf{Ende: 21.08 Uhr}

    % Wohl die komplizierteste Stelle im Protokoll
    % Als Faustregel sollten nicht mehr als 3 Namen auf einer Zeile stehen
    % Einfach an dieser Vorlage orientieren
    % Im Idealfall ist hier nicht viel zu ändern, da kaum jemand fehlt.
    % Rücktritte müssen natürlich pro Protokoll angepasst werden
    \vspace{0.5cm}
    \begin{tabular}{ll}
        Anwesend: & \\
            RCDS (5 Sitze): & Ramon Weilinger, Antonia Niecke, Ramin Shakiba, \\
                            & Jennifer Maack \\
             CampusGrün (14 Sitze): & Laura Franzen, Geoffrey Youett, Philipp Droll, \\
                                   & Tahnee Herzig, Svenja Horn, Jim Martens,\\
                                   & Mirzo Ulugbek Khatamov, Armin Günther, Martin Sievert\\
                                   & Elena Rysikova, Yasemin Günther \\
             Bier-Liste (2 Sitze) : & Jakob Pape, Karen Martirosian \\
             WiWi (2 Sitze): & Claas-Friso Hente \\
             Unicorns (5 Sitze): & Johannes Peplow, Annkathrin Löffler, \\
                                 & Andreas Hartkamp, Marielle Hermstrüwer \\
             Liste LINKS (3 Sitze): & Gunhild Berdal, Till Petersen, Sinah Mielich \\
             HWP (2 Sitze): & Ajdina Karahasan \\
             MIN (4 Sitze): & \\
             SDS* (3 Sitze): & Mena Winkler, Jacob Petersein, Artur Brückmann \\
             Bart-LISTE (2 Sitze): & Timo Zeimet, Dominic Laumer \\
             LHG (1 Sitz): & \\
             harte zeiten (1 Sitz): & Tobias Berking \\
             Jura (1 Sitz): & \\
             AL (2 Sitze): & Karima Schulze, Henri Weber \\
            & \\
        Entschuldigt: & Benjamin Welling (RCDS),\\
                      & Melf Johannsen, Mario Moldenhauer (CG), Samet Gunay,\\
                      & Kay Zöllmer (WiWi), Katharina Kucza (Unicorns), Ailina Salten, \\
                      & Lotte Rullkötter, Thea Wahlers (alle MIN), Tobias Heisig (LHG)\\
                      &\\
        Unentschuldigt abwesend: & Jan Detampel (MIN), \\
                                & Johann Baumhoefener (Jura)\\
                                &\\
        Rücktritte: & Freya Schmitz (CampusGrün) \rightarrow Sandra Fietz\\
    \end{tabular}
    \newpage
    %%% Vorgeschlagene Tagesordnung %%%
    \underline{Vorgeschlagene Tagesordnung}
    \begin{enumerate}[label={\textbf{Top \theenumi}},leftmargin=*]
        % folgende Zeile bei Sitzungsfortsetzungen einkommentieren und x durch
        % die Zahl des zuletzt behandelten TOPs ersetzen
        %\setcounter{enumi}{x}
        \item Keine Werbung von/für Bundeswehr (V1617-023) (60 Min.)
        \item Akkreditierung (V1617-024) (60 Min.)
        \item Verlegung der StuPa-Wahl in das Sommersemester (V1617-015) (60 Min.)
        \item Prüfauftrag: Wahl-O-Mat bei StuPa-Wahlen (V1617-022) (30 Min.)
        \item Verschiedenes
    \end{enumerate}

    \newpage

    % Hier beginnt Teil A. Bei Fortsetzungen kommt hier nur der Teil über
    % die Feststellung der endgültigen Fassung der Tagesordnung.
    % Für sonstige Sitzungen ist dies bereits das Gerüst für die Formalia.
    % Man muss nur noch die protokollierten Inhalte eintragen.

    {\Large \underline{Teil A}}

    \subsection{Formalia}

    \textit{Gunhild eröffnet Sitzung um 18.39 Uhr.}

    \subsubsection{Geschäftsbericht Präsidium}

    Es wird der Rücktritt von Freya Schmitz (CampusGrün) verlesen. Für Freya
    rückt Sandra Fietz nach.

    Der Besuch von der Wissenschaftssenatorin Katharina Fegebank erfolgt
    wahrscheinlich am 30. Juni für den Zeitraum von einer Stunde. Es wird davon
    ausgegangen, dass sie gegen 19 Uhr eintrifft.

    In Folge der Sitzungsverschiebung -- findet statt, weil am eigentlich
    für eine Sitzung geplanten 23. Juni bereits viele andere Veranstaltungen
    zeitgleich sind und der Haushalt spätestens am 23. Juni in eine erste
    Lesung im StuPa gehen muss -- finden die umgeplanten Sitzungen jetzt am
    16. und 30. Juni statt. Am 14. Juli bleibt wie zuvor eine Bedarfssitzung
    eingeplant.

    Auf der letzten Sitzung wurde das Präsidium damit beauftragt, eine Verfahrenslösung
    für den Umgang mit sog. Ersetzungsanträgen (Änderungsanträge, die den kompletten
    Text eines Antrags streichen und ersetzen) zu finden. Trotz einiger Unterhaltung
    hat das Präsidium noch keine abschließende Einigung gefunden.

    \subsubsection{Anfragen an das Präsidium}

    Es wird darauf hingewiesen, dass es am 9. Juni ein Freiraumfestival gibt,
    welches zeitgleich zur StuPa-Sitzung liegt. Daraufhin wird angemerkt,
    dass dies kein Problem sein müsse, da man die StuPa-Sitzung in den
    Festivaltag integrieren könne.

    Desweiteren wird angemerkt, dass StuPa-Protokolle durchaus kürzer gefasst
    und für eine längere Fassung auf die Videos verwiesen werden könnte.

    \subsubsection{Geschäftsbericht AStA}

    \textbf{Philipp Droll und Franziska Hildebrandt}: Es gab Gespräche mit
    Kanzler und weiteren Leuten. Dabei war mehr Anerkennung als Anecken vorhanden
    und es kam schnell zur Verständigung, sodass die Freiräume das Festival
    durchführen können. Auch der Haushaltsplan des Café Knallhart wurde besprochen.

    Der Raum 0029 im WiWi-Bunker könnte zum Gremienraum werden. Die Regelung für
    FSRe im Philturm (Anmerkung: bezieht sich auf den Zutritt zu den FSR-Räumen
    außerhalb der normalen Öffnungszeiten) könnte auf andere Gebäude übertragen
    werden. Es soll wohl schon sicher sein, dass diese Regelung auf den WiWi-Bunker
    ausgeweitet wird. Auch das Gebäude der Erziehungswissenschaften ist
    angedacht angedacht.

    Die generelle Nutzungsberechtigung für den Philturm wurde nur bis Ende 2016
    ausgegeben. Der Baubeginn für die Sanierung des Philturms könnte jedoch nicht
    direkt am Anfang von 2017 sein. Dies würde bedeuten, dass die Uni länger, als
    wegen des Bauhabens nötig, externe Räume anmieten muss. In jedem Fall wird es
    wohl Möglichkeiten geben Räume in den alten Postgebäuden anzumieten.

    In eineinhalb Wochen gibt es Gespräche mit dem Architekten der Uni.
    Der Vorplatz des Hauptgebäudes soll umgebaut werden. Die Schuldenuhr soll
    weiterhin stehenbleiben. Es gibt eine Diskussion über deren Ästhetik mit dem
    Kanzler.

    Die Kampagne gegen die Exzellenzinitiative geht weiter und viele
    Unterschriftenlisten aus den sind Fachbereichen zurückgekommen. Eine Einigung
    bis zum 16.6. bezüglich der Exzellenzinitiative ist unwahrscheinlich. Nächste
    Woche findet am 2.6. eine studentische Vollversammlung zum Thema Exzellenzinitiative
    statt.

    In der letzten Vorlesungswoche soll es eine Konferenz zu HealthEconomics geben.
    Die Veranstaltungen im ESA-Flügel sollen dafür ausfallen. Der ermäßigte Preis
    für Studierende beträgt 300€.

    Eine Veranstaltung mit Frau Leopold kommt nicht zustande, weil sie (angeblich)
    keine Zeit hat.

    Die KMK\footnote{Kultusministerkonferenz} denkt über eine BA/MA-Reform nach.
    Das erste Jahr im Studium soll danach notenfrei werden.

    \textbf{Timo Tomanek?} berichtet über die Finanzen. Nächste Woche soll der
    erste Entwurf des Haushalts AStA-intern vorgelegt werden.

    \textbf{Geoffrey Youett} hält den Bericht aus dem Referat
    ÖffKulturHoPo\footnote{Öffentlichkeitsarbeit, Kultur, Hochschulpolitik}. Es
    wurde ein Postverteiler aufgebaut, über den allen FSRen alle 2 Wochen Material
    zugeschickt werden soll. Es werden sich darüber Gedanken gemacht, wie die FSRe
    besser in die inhaltliche Arbeit und Demos einbezogen werden können.

    Außerdem soll es eine Fotoaktion zur Frage "Was bedeutet Uni für dich?" geben.
    Ebenfalls wurden sich darüber Gedanken gemacht, welche Veranstaltungen in Bezug
    zur EM angeboten werden könnten.

    Die Frage der Gestaltung des Campus OpenAir ist ein Dauerbrenner bei der
    AStA-Arbeit. Dabei lag und liegt der Hauptstreitpunkt bei der Frage, ob das
    Campus OpenAir ein studentisch selbstorganisiertes Festival sein soll oder ob
    Externe in die Organisation und Durchführung eingebunden werden sollen.
    Das Freiraum-Festival ist im Prinzip ein Campus OpenAir, der ersteren Art:
    Es ist ein studentisch organisiertes Festival ohne externe Firma.

    Der AStA steht als Kooperationspartner auf den Plakaten drauf. Die Organisation
    und Planung wurde jedoch alleine von den Freiräumen durchgeführt. Das Festival
    beginnt um 16 Uhr. Um 18 Uhr wird die Musik lauter gedreht bis 22 Uhr. Eventuell
    wird es eine Sondergenehmigung geben, um bis Mitternacht (0 Uhr) Musik machen
    zu können.

    Es wird auch Stände auf dem Festival geben. Der AStA hat sich überlegt, wie
    er sich dort einbringen kann. Eine Idee ist es, dass sich der AStA mit
     EM-Veranstaltungen einbringt.

    \textit{Zeitansage: 19.14 Uhr}

    Die Kulturkurse laufen alle. Es gibt jedoch einige Kinderkrankheiten. Das
    Buchungssystem soll optimiert und eine neue Website gestaltet werden. Im
    nächsten Semester soll es noch besser und fruchtbarer ablaufen.

    \textbf{Jacob Petersein} berichtet aus dem Referat für Antidiskriminierung,
    Anfifaschismus und Kritische Wissenschaft. Das Referat kann bislang nicht
    wirklich mit der Arbeit beginnen, da viele Personen nicht da sind.

    Das übergreifende Programm des AStA sind die EM-Veranstaltungen. Es ist das
    Ziel die Länder der gegnerischen Mannschaften der deutschen Nationalmannschaft
    vorzustellen und parallel zu den Spielen ein OpenAir-Kino anzubieten. Außerdem
    gibt es Überlegungen zu weiteren Veranstaltungsarten. Am 9.7. soll es ein
    Fußballturnier in Kooperation mit Refugees Welcome Karoviertel stattfinden.

    Die Veranstaltung zur Rolle des NS-Studentenbundes an der Uni Hamburg war
    gut besucht und hat eine heterogene Zuschauermenge angezogen. Jacob fasst
    die Veranstaltung weitergehend zusammen.

    \textbf{Artur Brückmann} berichtet vom AS\footnote{Akademischer Senat}.
    Die Vertreter für den Hochschulrat wurden gewählt: Herr Nöth und Frau Platzer
    haben jeweils die 10 Stimmen der Professor*Innen auf sich vereinen können.
    Bei den Ziel- und Leistungsvereinbarungen haben sich die Unileitung und die
    Behörde darauf geeinigt die quantitativen Kriterien niedrig anzusetzen, sodass
    diese tatsächlich erfüllbar sind. Dies sei der erste Schritt auf dem Weg zu
    qualitativen Kriterien. Schließlich beschließt der AS mit 16:0:2, dass der
    Aufruf der Wissenschaftler*Innen für Frieden unterzeichnet wird.

    Mit dem Studierendenwerk wurde vereinbart, dass die
    Wissenschaftsausschuss-Sitzung gemeinsam vorbereitet wird.

    \textbf{Armin Günther} erwähnt, dass am 30.5. im ESA B die letzte Veranstaltung
    der Reihe "ich konsumiere, du konsumierst, wir kritisieren" stattfinden wird.

    Außerdem wird es eine Veranstaltung zum Thema Exzellenzinitiative geben --
    zusammen mit WiMis der Uni Jena.

    \textbf{Oliver Vornfeld} berichtet aus dem Ausschuss gegen Rechts. Am 5.7.
    wird es den Tag gegen Rechts geben.

    \subsubsection{Anfragen an den AStA}

    Till Petersen interessiert sich dafür, was für ein Termin die Ausschusssitzung
    mit dem Studierendenwerk ist.

    Franziska erklärt, wie es zur Sitzung gekommen ist. Artur ergänzt.

    \textit{Zeitansage: 19.30 Uhr}

    Vor dem Hintergrund des Designgrundsatzes "Form follows Function" fragt Till,
    inwieweit das Öffentlichkeitsreferat die Form bzgl. des zu verbreitenden
    Inhalts reflektiert hat.

    Es wird von Geoffrey geantwortet. Ergänzend wird gesagt, dass der Fokus
    bei Veranstaltungen auf dem persönlichen Kontakt liegen soll.

    Till erwähnt, dass nicht alle FSRe mit der Regelung im Philturm zufrieden
    sind. Von daher sei die Verallgemeinerung dessen nicht per se sinnvoll. Für
    eine Ausdehnung auf das Gebäude der Erziehungswissenschaften müssten zum
    Beispiel die FSRe der Erziehungswissenschaften eingebunden werden.

    Franziska findet die Regelung im Philturm auch nicht ideal. Die Regelung sei
    ein Zwischenstand. Zwischen den FSRen verliefe eine Konfliktlinie, wofür längerer
    Zugang benötigt wird. Die einen wollten arbeiten, die anderen Partys feiern.
    In Bezug zu VMP 8 gebe es Gespräche.

    Geoffrey wird gefragt, was mit dem Begriff Freiraumkultur gemeint ist.

    Er entgegnet, dass das Referat erst spät in die Diskussion um das Festival
    eingestiegen sei. Bei den vorbereitenden Gesprächen habe es Einigkeit darüber
    gegeben, dass die Existenz der Freiräume auf einem politischen Hintergrund beruht.
    Man habe sich in Bezug zum Festival vorgenommen, diese Diskussion über die
    Bedeutung des Wortes Freiraumkultur in den folgenden Monaten fortzuführen.

    \textit{Zeitansage: 19.47 Uhr}

    \textit{Feststellung der Beschlussfähigkeit: Mit 25/47 Parlamentarier*Innen
    ist diese gegeben.}

    \subsubsection{Dringlichkeitsanträge des AStA}

    Keine.

    \subsubsection{Aktuelle Stunde (falls entsprechender Antrag vorliegt)}

    Es wurde kein Antrag gestellt.

    \subsubsection{Feststellung der endgültigen Fassung des Teils B der Tagesordnung}

    Jacob möchte TOP 4 und 3 tauschen. Die Mehrheit ist dafür (15:3:4).

    Die so geänderte Tagesordnung wurde ohne Widerspruch angenommen.

    %%% BESCHLOSSENE TAGESORDNUNG %%%
    \underline{Endgültige Fassung der Tagesordnung}
    \begin{enumerate}[label={\textbf{Top \theenumi}},leftmargin=*]
        % folgende Zeile bei Sitzungsfortsetzungen einkommentieren und x durch
        % die Zahl des zuletzt behandelten TOPs ersetzen
        %\setcounter{enumi}{x}
        \item Keine Werbung von/für Bundeswehr (V1617-023) (60 Min.)
        \item Akkreditierung (V1617-024) (60 Min.)
        \item Prüfauftrag: Wahl-O-Mat bei StuPa-Wahlen (V1617-022) (30 Min.)
        \item Verlegung der StuPa-Wahl in das Sommersemester (V1617-015) (60 Min.)
        \item Verschiedenes
    \end{enumerate}

    \subsubsection{Feststellung der Beschlussfähigkeit}

    Schon geschehen.

    \subsubsection{Genehmigung der Protokolle der vorangegangenen Sitzungen}

    Das Protokoll der 1. Fortsetzung der ersten Sitzung wird bei einer Enthaltung
    und 0 Gegenstimmen mehrheitlich angenommen.

    Das Protokoll der 2. Sitzung  wird bei 3 Enthaltungen und 0 Gegenstimmen
    angenommen.

    \textit{Zeitansage: 20.00 Uhr}

    \textit{Präsidiumspause bis 20.10 Uhr (angekündigt)/20.15 Uhr (defacto)}

    \newpage
    %\vspace{0.5cm}
    {\Large \underline{Teil B}}

    % Teil B beginnt
    % Dieser wird sich für jede Sitzung unterscheiden.
    % Jeder TOP ist eine subsection.
    % Vorhandene Beispiele anfangs entfernen und durch tatsächliche TOPs ersetzen

    % Folgende Zeile bei Sitzungsfortsetzungen einkommentieren und x mit
    % der Zahl des zuletzt behandelten TOPs ersetzen.
    %\setcounter{subsection}{x}
    \subsection{Keine Werbung von/für Bundeswehr}

    Franzi stellt den Antrag vor.

    Tobias möchte, dass die beiden Absätze des Antrags getrennt werden
    und eigenständig behandelt werden. Außerdem hat er einen Änderungsantrag:

    Zeile 4 von Universität bis Zeile 5 werden ersetzen durch:
    Wissenschaft hat die Aufgabe, die Wahrheit für die menschenwürdige
    Entwicklung der Welt zu bestimmen und zur Wirklichkeit zu machen. Die
    Wahrheit stirbt jedoch im Krieg bekanntlich zuerst.

    Artur hält einen Redebeitrag und erklärt den historischen Zusammenhang
    der Gründung der Uni Marketing GmbH.

    20:32

    Gunhild möchte "künftigen" in Zeile 7 streichen. Am 11. Juni findet
    der Tag der Bundeswehr statt. Gunhild befürwortet die Teilung des Antrags.

    Ramin stellt GO-Antrag auf Feststellg. der Beschlussfähigkeit. Till beantragt
    eine Fraktionspause für Liste Links von 20 Minuten bis 21:03.

    In dieser Zeit wird weiter diskutiert.

    Geoffrey hat eine persönliche Erklärung abgegeben:

    Sinnbildlich spiegeln die gähnend leeren Sitzplätze der Opposition
    wider, was hinter ihren Wahlprogrammen steckt -- nichts.

    Jetzt folgt die Abstimmung: Mit 22 Parlamentarier*Innen nicht beschlussfähig.

    Sitzung ist abgebrochen um 21:08. Nächste Sitzung am 9. Juni.
\end{document}
